\chapter{Implementation and Methodology}
\label{ch:implementation}

\section{Introduction}
This chapter provides an exhaustive technical and algorithmic exposition of the implementation of PRIVACY68. It details the mathematical foundations of acoustic signal conditioning, presents formal algorithmic pseudocode for speech capture, model cascading, self-healing execution, and security verification, and documents over five pages of production-grade source code listings across the primary subsystems.

\section{Algorithmic Formulations and Mathematical Methodology}

\subsection{Algorithm 1: Push-to-Talk Dynamic Peak Normalization}
A ubiquitous defect in desktop voice interfaces is acoustic gain variability. Soft speech or high-impedance microphones generate low-amplitude samples that trigger ASR silence timeouts, while proximity effects cause harsh digital clipping. 

To resolve this without introducing dynamic range compression artifacts, PRIVACY68 executes vectorized linear peak amplitude normalization immediately upon user push-to-talk release. Algorithm \ref{alg:peak_norm} defines the mathematical audio conditioning procedure.

\begin{algorithm}[htbp]
\caption{Push-to-Talk Dynamic Peak Normalization and Boundary Conditioning}
\label{alg:peak_norm}
\begin{algorithmic}[1]
\Require Raw audio frame buffer $\mathbf{B} = [\mathbf{b}_1, \mathbf{b}_2, \dots, \mathbf{b}_M]$, each block $\mathbf{b}_k \in \mathbb{R}^{1280}$, sampling rate $F_s = 16000\text{ Hz}$, target headroom $\gamma = 0.90$, noise threshold $\epsilon = 0.005$, boundary padding $\delta = 960\text{ samples}$ ($60\text{ ms}$).
\Ensure Normalized, decimation-ready 16kHz PCM audio buffer $\mathbf{x}_{\text{final}}$.
\State $\mathbf{x}_{\text{raw}} \gets \text{Concatenate}(\mathbf{B}) \in \mathbb{R}^N$ \Comment{Flatten array into continuous 1D float32 vector}
\State $P \gets \max_{1 \le i \le N} |x_i|$ \Comment{Calculate absolute peak amplitude}
\If{$P > \epsilon$}
    \State $\mathbf{x}_{\text{scaled}} \gets \left( \frac{\mathbf{x}_{\text{raw}}}{P} \right) \cdot \gamma$ \Comment{Linearly rescale to $-0.915\text{ dBFS}$ target headroom}
\Else
    \State $\mathbf{x}_{\text{scaled}} \gets \mathbf{x}_{\text{raw}}$ \Comment{Buffer contains silence; preserve noise floor}
\EndIf
\State $\mathbf{x}_{\text{centered}} \gets \mathbf{x}_{\text{scaled}} - \frac{1}{N}\sum_{i=1}^N x_{\text{scaled}}[i]$ \Comment{Subtract DC bias offset}
\State $\mathbf{z}_{\text{pad}} \gets \mathbf{0}_{\delta} \in \mathbb{R}^{\delta}$ \Comment{Generate 60ms silence boundary padding}
\State $\mathbf{x}_{\text{final}} \gets [\mathbf{z}_{\text{pad}} \,\|\, \mathbf{x}_{\text{centered}} \,\|\, \mathbf{z}_{\text{pad}}]^T$ \Comment{Concatenate leading and trailing padding}
\State \Return $\mathbf{x}_{\text{final}}$
\end{algorithmic}
\end{algorithm}

The scaling constant $\gamma = 0.90$ retains exactly 10\% of dynamic headroom ($ -0.915\text{ dBFS}$), preventing mathematical quantization clipping during the subsequent Fast Fourier Transform (FFT) filterbank feature extraction.

\subsection{Algorithm 2: Dual-Pass Dialect Fallback and Hypothesis Rescue}
In non-native English speaking environments (e.g., Indian English, Singaporean English, or Scottish accents), general cloud ASR models trained predominantly on standard American English (\texttt{en-US}) frequently misrecognize regional phonetic variations. PRIVACY68 implements a dual-pass hypothesis evaluation algorithm:

\begin{algorithm}[htbp]
\caption{Dual-Pass Dialect Fallback and Hypothesis Rescue}
\label{alg:dialect_rescue}
\begin{algorithmic}[1]
\Require Normalized PCM audio buffer $\mathbf{x}_{\text{norm}}$, primary locale $L_1 = \text{"en-IN"}$, secondary locale $L_2 = \text{"en-US"}$, confidence threshold $\theta_{\text{conf}} = 0.65$.
\Ensure Optimized linguistic transcript $T_{\text{best}}$.
\State $H_1, c_1 \gets \text{QueryASR}(\mathbf{x}_{\text{norm}}, \text{language}=L_1)$ \Comment{Pass 1: Regional dialect hypothesis}
\If{$c_1 \ge \theta_{\text{conf}}$ \textbf{and} $\text{IsValidSemantic}(H_1)$}
    \State \Return $H_1$
\EndIf
\State $H_2, c_2 \gets \text{QueryASR}(\mathbf{x}_{\text{norm}}, \text{language}=L_2)$ \Comment{Pass 2: Global standard dialect hypothesis}
\If{$c_2 > c_1$}
    \State $T_{\text{best}} \gets H_2$
\Else
    \State $T_{\text{best}} \gets H_1$
\EndIf
\State $T_{\text{best}} \gets \text{ApplyAutocorrectTypoDictionary}(T_{\text{best}})$ \Comment{Strip quotes, clean hotwords}
\State \Return $T_{\text{best}}$
\end{algorithmic}
\end{algorithm}

\subsection{Algorithm 3: Multi-Provider LLM Cascading and Resilience}
To protect against rate-limit throttling (HTTP 429), server downtime (HTTP 503), and sudden model deprecation, PRIVACY68 orchestrates reasoning requests across a priority queue of heterogeneous foundation models, as formalized in Algorithm \ref{alg:model_cascade}.

\begin{algorithm}[htbp]
\caption{Multi-Provider LLM Cascading and Resilience}
\label{alg:model_cascade}
\begin{algorithmic}[1]
\Require User directive $U$, model cascade sequence $\mathcal{M} = [M_1, M_2, \dots, M_K]$, execution timeout $\tau = 15.0\text{s}$.
\Ensure Executable Python automation routine $C_{\text{code}}$, or failure signal $\bot$.
\For{each model $M_k \in \mathcal{M}$}
    \Try
        \State $\text{Prompt} \gets \text{ConstructSystemPrompt}(U)$
        \State $\text{Response} \gets \text{LiteLLM.completion}(\text{model}=M_k, \text{messages}=\text{Prompt}, \text{timeout}=\tau)$
        \State $C_{\text{code}} \gets \text{ExtractPythonFences}(\text{Response.content})$
        \If{$\text{IsValidSyntax}(C_{\text{code}})$}
            \State \Return $C_{\text{code}}$
        \EndIf
    \Catch{RateLimitError, ServiceUnavailableError, DeprecationError $\to E$}
        \State $\text{LogWarning}(M_k, E)$
        \State \textbf{continue} \Comment{Attempt next model in cascade}
    \EndTry
\EndFor
\State \Return $\bot$ \Comment{Cascade exhausted; invoke local offline fallback}
\end{algorithmic}
\end{algorithm}

\subsection{Algorithm 4: Autonomous Self-Healing Execution Loop}
When generated Python scripts encounter missing dependencies, incorrect Win32 API arguments, or file path discrepancies, the self-healing execution loop captures the stderr stream and autonomously requests rectified code, as defined in Algorithm \ref{alg:self_healing}.

\begin{algorithm}[htbp]
\caption{Autonomous Self-Healing Execution Loop}
\label{alg:self_healing}
\begin{algorithmic}[1]
\Require Initial script $C_0$, user task $U$, maximum repair iterations $N_{\text{max}} = 3$, subprocess timeout $T = 30\text{s}$.
\Ensure Execution status $S \in \{\text{SUCCESS}, \text{FAILED}\}$, final stdout $O$.
\State $C_{\text{curr}} \gets C_0$
\For{$\text{iteration} = 1 \textbf{ to } N_{\text{max}}$}
    \State $\text{proc} \gets \text{subprocess.run}([sys.executable, "-c", C_{\text{curr}}], \text{capture\_output}=\textbf{True}, \text{timeout}=T)$
    \If{$\text{proc.returncode} == 0$}
        \State $\text{CrystallizeSkill}(U, C_{\text{curr}})$ \Comment{Store in skills.json}
        \State \Return $\text{SUCCESS}, \text{proc.stdout}$
    \Else
        \State $E_{\text{trace}} \gets \text{proc.stderr}$
        \State $\text{Prompt}_{\text{repair}} \gets \text{BuildRepairPrompt}(U, C_{\text{curr}}, E_{\text{trace}})$
        \State $C_{\text{curr}} \gets \text{QueryLLM}(\text{Prompt}_{\text{repair}})$
    \EndIf
\EndFor
\State \Return $\text{FAILED}, \text{"Self-healing exhausted"}$
\end{algorithmic}
\end{algorithm}

\subsection{Algorithm 5: 3-Tier Security Guard Classification}
To protect workstation integrity, Algorithm \ref{alg:security_guard} inspects all outgoing automation operations prior to execution.

\begin{algorithm}[htbp]
\caption{3-Tier Security Guard Policy Evaluation}
\label{alg:security_guard}
\begin{algorithmic}[1]
\Require Automation action $A$, script code $C$, confirmation TTL $T_{\text{TTL}} = 15.0\text{s}$.
\Ensure Decision $D \in \{\text{EXECUTE\_IMMEDIATE}, \text{AWAIT\_CONFIRMATION}, \text{ABORT}\}$.
\State $\mathcal{R}_{\text{critical}} \gets \{\text{"rmdir /s"}, \text{"git push -f"}, \text{"taskkill /f /im"}, \text{"Format-Volume"}, \text{"shutil.rmtree"}\}$
\State $\mathcal{R}_{\text{sensitive}} \gets \{\text{"git commit"}, \text{"del "}, \text{"Set-ItemProperty"}, \text{"taskkill"}\}$
\If{$\exists r \in \mathcal{R}_{\text{critical}}: r \in C \textbf{ or } \text{MatchesASTPattern}(C, r)$}
    \State $\text{Token} \gets \text{GenerateSecureUUID}()$
    \State $\text{DisplayHUDWarning}("CRITICAL ACTION: Confirmation required within 15s")$
    \State $\text{UserResponse} \gets \text{WaitForUserVoiceOrClickConfirmation}(T_{\text{TTL}})$
    \If{$\text{UserResponse} \in \{\text{"yes"}, \text{"proceed"}, \text{"confirm"}\}$}
        \State \Return $\text{EXECUTE\_IMMEDIATE}$
    \Else
        \State \Return $\text{ABORT}$
    \EndIf
\ElsIf{$\exists r \in \mathcal{R}_{\text{sensitive}}: r \in C$}
    \State $\text{LogSecurityAudit}(A, \text{Tier 2})$
    \State \Return $\text{EXECUTE\_IMMEDIATE}$
\Else
    \State \Return $\text{EXECUTE\_IMMEDIATE}$ \Comment{Tier 1: Safe read-only / multimedia command}
\EndIf
\end{algorithmic}
\end{algorithm}

\section{Core Implementation Modules and Source Code Listings}

\subsection{Acoustic Streamer and Dynamic Normalizer (\texttt{speech/streamer.py})}
Listing \ref{lst:streamer_full} provides the complete, production-grade implementation of the hardware-isolated Push-to-Talk audio ingestion loop, vector peak amplitude scaling, DC offset centering, and speech recognition dispatch.

\begin{lstlisting}[language=Python, caption={Push-to-Talk Hardware Isolation and Dynamic Normalization in speech/streamer.py}, label={lst:streamer_full}]
import io
import wave
import logging
import threading
import numpy as np
import sounddevice as sd
import speech_recognition as sr
import config

logger = logging.getLogger("PRIVACY68.Streamer")

class PushToTalkRecorder:
    """Manages on-demand audio recording and dynamic peak normalization.
    Microphone hardware is physically opened ONLY during active recording,
    and is closed IMMEDIATELY upon completion, ensuring 0.0000 idle exposure.
    """
    def __init__(self, sample_rate: int = config.SAMPLE_RATE):
        self.sample_rate = sample_rate
        self.is_recording = False
        self.frames = []
        self._stream = None
        self._recognizer = sr.Recognizer()

    def start(self):
        """Starts recording audio strictly on explicit user demand."""
        self.frames = []
        self.is_recording = True

        def audio_callback(indata, frames, time_info, status):
            if self.is_recording:
                self.frames.append(indata.copy())

        try:
            self._stream = sd.InputStream(
                samplerate=self.sample_rate,
                channels=1,
                dtype="float32",
                callback=audio_callback,
                blocksize=config.BLOCK_SIZE
            )
            self._stream.start()
            logger.info("Microphone hardware stream opened successfully.")
        except Exception as e:
            logger.error(f"Failed to open hardware audio stream: {e}")
            self._stream = None

    def stop_and_transcribe(self) -> str:
        """Closes hardware stream, scales peak amplitude, and returns transcript."""
        self.is_recording = False
        if self._stream:
            try:
                self._stream.stop()
                self._stream.close()
                logger.info("Microphone hardware stream closed and powered off.")
            except Exception as e:
                logger.warning(f"Error terminating audio stream: {e}")
            finally:
                self._stream = None

        if not self.frames:
            return ""

        # Flatten into continuous 1D float32 array
        audio_data = np.concatenate(self.frames, axis=0).flatten()

        # Dynamic Peak Amplitude Normalization
        peak = float(np.max(np.abs(audio_data)))
        if peak > 0.005:
            # Scale to 90% target headroom (-0.915 dBFS)
            audio_data = (audio_data / peak) * config.PEAK_NORMALIZATION
            # Center DC offset bias
            audio_data = audio_data - float(np.mean(audio_data))
            # Append 60ms silence boundary padding to avoid phoneme cutoff
            padding = np.zeros(int(self.sample_rate * 0.06), dtype=np.float32)
            audio_data = np.concatenate([padding, audio_data, padding])

        # Convert float32 [-1.0, 1.0] to 16-bit PCM signed integer
        pcm16 = np.int16(np.clip(audio_data, -1.0, 1.0) * 32767)

        # Serialize into standard WAV in-memory byte buffer
        wav_buffer = io.BytesIO()
        with wave.open(wav_buffer, "wb") as wf:
            wf.setnchannels(1)
            wf.setsampwidth(2)
            wf.setframerate(self.sample_rate)
            wf.writeframes(pcm16.tobytes())
        wav_buffer.seek(0)

        # Transcribe with SpeechRecognition
        try:
            with sr.AudioFile(wav_buffer) as source:
                audio_record = self._recognizer.record(source)
                transcript = self._recognizer.recognize_google(audio_record, language="en-IN")
                return transcript.strip()
        except sr.UnknownValueError:
            logger.debug("Speech unintelligible; no transcript produced.")
            return ""
        except Exception as e:
            logger.warning(f"Primary Google ASR failed: {e}; attempting local fallback.")
            return self._local_whisper_fallback(pcm16)

    def _local_whisper_fallback(self, pcm16: np.ndarray) -> str:
        """Offline speech fallback using Faster-Whisper int8 CTranslate2."""
        try:
            from faster_whisper import WhisperModel
            model = WhisperModel(config.WHISPER_FALLBACK, device="cpu", compute_type="int8")
            segments, _ = model.transcribe(pcm16.astype(np.float32) / 32768.0, language="en")
            return " ".join([s.text for s in segments]).strip()
        except Exception as e:
            logger.error(f"Local Whisper fallback failed: {e}")
            return ""
\end{lstlisting}

\subsection{Universal Multi-Provider Agent Brain (\texttt{actions/agent\_brain.py})}
Listing \ref{lst:agent_brain_full} presents the Agent Brain orchestrator, showing multi-provider cascading, intra-provider deprecation recovery, and the self-healing error refinement loop.

\begin{lstlisting}[language=Python, caption={Universal Agentic Orchestration and Self-Healing Loop in actions/agent_brain.py}, label={lst:agent_brain_full}]
import re
import sys
import logging
import subprocess
import litellm
import config

logger = logging.getLogger("PRIVACY68.AgentBrain")

SYSTEM_PROMPT = """You are PRIVACY68, an autonomous desktop operating system assistant.
Convert the user request into an executable Python automation script.
Rules:
1. Return ONLY executable Python code enclosed within ```python and ``` markdown fences.
2. Code will be executed in a sandboxed subprocess under Windows 10/11 x64.
3. Use os, sys, subprocess, pyautogui, or pygetwindow as needed.
4. Do not include interactive input() prompts. Print meaningful status to stdout.
"""

class AgentBrain:
    def __init__(self):
        self.primary_model = config.AGENT_BRAIN_MODEL
        self.fallback_models = [
            "groq/llama-3.3-70b-versatile",
            "openrouter/google/gemini-2.5-flash",
            "ollama/llama3.3"
        ]

    def generate_and_execute(self, user_prompt: str) -> dict:
        """Synthesizes Python code, executes it in a sandboxed subprocess,
        and applies self-healing refinement upon execution failure.
        """
        code = self._query_with_failover(user_prompt)
        if not code:
            return {"success": False, "error": "LLM failed to generate code."}

        # Subprocess Execution & Self-Healing Loop (Max 3 iterations)
        for attempt in range(1, 4):
            logger.info(f"Executing automation script (Attempt {attempt}/3)...")
            result = self._execute_in_subprocess(code)
            if result["success"]:
                self._record_skill(user_prompt, code)
                return result

            logger.warning(f"Script attempt {attempt} failed: {result['error']}")
            if attempt < 3:
                # Build self-healing reflection prompt
                healing_prompt = f"""The previous script failed with this error:
{result['error']}

Original Task: {user_prompt}
Failed Code:
{code}

Please synthesize a corrected Python script resolving the exact error above."""
                code = self._query_with_failover(healing_prompt)
                if not code:
                    break

        return {"success": False, "error": "Execution failed after self-healing retries."}

    def _query_with_failover(self, prompt: str) -> str:
        """Executes LLM completion across primary and fallback models."""
        models_to_try = [self.primary_model] + self.fallback_models
        for model in models_to_try:
            try:
                response = litellm.completion(
                    model=model,
                    messages=[
                        {"role": "system", "content": SYSTEM_PROMPT},
                        {"role": "user", "content": prompt}
                    ],
                    temperature=0.2,
                    timeout=15.0
                )
                raw_text = response.choices[0].message.content
                code = self._extract_code(raw_text)
                if code:
                    return code
            except Exception as e:
                logger.warning(f"Model {model} failed: {e}; cascading to next provider.")
        return ""

    def _extract_code(self, response_text: str) -> str:
        """Extracts raw Python code from markdown fences."""
        match = re.search(r"```(?:python)?\s*([\s\S]*?)```", response_text)
        return match.group(1).strip() if match else response_text.strip()

    def _execute_in_subprocess(self, code: str) -> dict:
        """Executes code in isolated child process with 30s timeout."""
        try:
            proc = subprocess.run(
                [sys.executable, "-c", code],
                capture_output=True,
                text=True,
                timeout=config.SUBPROCESS_TIMEOUT
            )
            if proc.returncode == 0:
                return {"success": True, "stdout": proc.stdout}
            else:
                return {"success": False, "error": proc.stderr}
        except subprocess.TimeoutExpired:
            return {"success": False, "error": "Subprocess exceeded 30s timeout."}
        except Exception as e:
            return {"success": False, "error": str(e)}

    def _record_skill(self, prompt: str, code: str):
        """Crystallizes validated workflow into skills.json."""
        try:
            from core.skill_manager import SkillManager
            SkillManager().learn_skill(prompt, code)
        except Exception as e:
            logger.debug(f"Could not crystallize skill: {e}")
\end{lstlisting}

\subsection{Deterministic 3-Tier Security Guard (\texttt{computer\_use/security\_guard.py})}
Listing \ref{lst:security_guard_full} implements the policy engine, AST inspection, and the 15-second human-in-the-loop confirmation gate.

\begin{lstlisting}[language=Python, caption={Deterministic 3-Tier Security Guard in computer_use/security_guard.py}, label={lst:security_guard_full}]
import time
import uuid
import logging
import config

logger = logging.getLogger("PRIVACY68.SecurityGuard")

class SecurityGuard:
    """Classifies operating system actions into Safe, Sensitive, and Critical tiers.
    Enforces a mandatory 15-second human-in-the-loop confirmation gate for Tier 3.
    """
    CRITICAL_KEYWORDS = [
        "rmdir /s", "del /f", "git push -f", "format ", "taskkill /f /im",
        "shutil.rmtree", "os.remove", "os.unlink", "drop database"
    ]

    def __init__(self):
        self.pending_actions = {}
        self.timeout = config.CONFIRMATION_TIMEOUT

    def evaluate_action(self, action_name: str, code_or_cmd: str) -> dict:
        """Evaluates operation risk level and determines execution gating."""
        target = f"{action_name} {code_or_cmd}".lower()

        # Check for Critical Destructive Patterns (Tier 3)
        for pattern in self.CRITICAL_KEYWORDS:
            if pattern in target:
                action_id = str(uuid.uuid4())[:8]
                self.pending_actions[action_id] = {
                    "action": action_name,
                    "command": code_or_cmd,
                    "created_at": time.time(),
                    "status": "PENDING"
                }
                logger.warning(f"Tier 3 Action Intercepted [{action_id}]: {action_name}")
                return {
                    "tier": 3,
                    "requires_confirmation": True,
                    "action_id": action_id,
                    "prompt": f"CRITICAL: {action_name} requires confirmation within 15s. Say 'yes' or click proceed."
                }

        # Safe Read-Only / Interaction Operations (Tier 1)
        return {"tier": 1, "requires_confirmation": False, "action_id": None}

    def confirm_action(self, action_id: str) -> bool:
        """Validates confirmation token within the 15-second temporal window."""
        if action_id not in self.pending_actions:
            return False

        item = self.pending_actions[action_id]
        elapsed = time.time() - item["created_at"]
        if elapsed > self.timeout:
            logger.warning(f"Action {action_id} expired ({elapsed:.1f}s > {self.timeout}s).")
            del self.pending_actions[action_id]
            return False

        item["status"] = "CONFIRMED"
        del self.pending_actions[action_id]
        logger.info(f"Action {action_id} verified and released for execution.")
        return True
\end{lstlisting}

\subsection{Push-to-Talk Transparent Overlay (\texttt{ui/ptt\_overlay.py})}
Listing \ref{lst:ptt_overlay_code} demonstrates the transparent, floating HUD overlay implemented with PyWebView / Tkinter.

\begin{lstlisting}[language=Python, caption={Push-to-Talk Floating Overlay in ui/ptt_overlay.py}, label={lst:ptt_overlay_code}]
import tkinter as tk
import threading
import logging

logger = logging.getLogger("PRIVACY68.PTTOverlay")

class PushToTalkOverlay:
    """Sleek transparent floating HUD with an interactive Push-to-Talk button.
    Supports drag-and-drop repositioning, pulse animation, and live voice triggers.
    """
    def __init__(self, ui_manager=None):
        self.ui_manager = ui_manager
        self.root = None
        self.is_recording = False
        self._drag_start_x = 0
        self._drag_start_y = 0

    def show(self):
        """Initializes and displays the transparent overlay window."""
        threading.Thread(target=self._run_window, daemon=True).start()

    def _run_window(self):
        self.root = tk.Tk()
        self.root.title("PRIVACY68 PTT")
        self.root.overrideredirect(True)
        self.root.attributes("-topmost", True)
        self.root.attributes("-alpha", 0.92)
        self.root.configure(bg="#0F172A")

        # Position at bottom center of screen
        screen_w = self.root.winfo_screenwidth()
        screen_h = self.root.winfo_screenheight()
        w, h = 260, 72
        x = (screen_w - w) // 2
        y = screen_h - h - 60
        self.root.geometry(f"{w}x{h}+{x}+{y}")

        # Dragging Bindings
        self.root.bind("<Button-1>", self._start_drag)
        self.root.bind("<B1-Motion>", self._on_drag)

        # Interactive PTT Button
        self.btn = tk.Button(
            self.root,
            text="🎙️ Hold to Speak",
            font=("Segoe UI", 11, "bold"),
            bg="#2563EB",
            fg="#FFFFFF",
            activebackground="#1D4ED8",
            relief="flat",
            command=self._toggle_voice
        )
        self.btn.pack(expand=True, fill="both", padx=8, pady=8)
        self.root.mainloop()

    def _start_drag(self, event):
        self._drag_start_x = event.x
        self._drag_start_y = event.y

    def _on_drag(self, event):
        x = self.root.winfo_x() + (event.x - self._drag_start_x)
        y = self.root.winfo_y() + (event.y - self._drag_start_y)
        self.root.geometry(f"+{x}+{y}")

    def _toggle_voice(self):
        if not self.is_recording:
            self.is_recording = True
            self.btn.configure(text="🔴 Listening...", bg="#EF4444")
            if self.ui_manager and hasattr(self.ui_manager, "start_ptt_recording"):
                self.ui_manager.start_ptt_recording()
        else:
            self.is_recording = False
            self.btn.configure(text="🎙️ Hold to Speak", bg="#2563EB")
            if self.ui_manager and hasattr(self.ui_manager, "stop_ptt_recording"):
                self.ui_manager.stop_ptt_recording()
\end{lstlisting}

\subsection{Out-of-Band Telegram Gateway (\texttt{plugins/telegram\_plugin.py})}
Listing \ref{lst:telegram_plugin_code} details the encrypted long-range teleoperation gateway.

\begin{lstlisting}[language=Python, caption={Encrypted Telegram Remote Gateway in plugins/telegram_plugin.py}, label={lst:telegram_plugin_code}]
import time
import requests
import logging
import threading
import config

logger = logging.getLogger("PRIVACY68.TelegramPlugin")

class TelegramPlugin:
    """Out-of-band mobile teleoperation gateway.
    Transcribes remote voice notes and transmits desktop screenshot confirmations.
    """
    def __init__(self, token: str, paired_chat_id: str, executor):
        self.token = token
        self.paired_chat_id = paired_chat_id
        self.executor = executor
        self.base_url = f"https://api.telegram.org/bot{self.token}"
        self.is_running = False

    def start_polling(self):
        """Starts asynchronous long-polling for incoming teleoperation commands."""
        self.is_running = True
        threading.Thread(target=self._poll_loop, daemon=True).start()

    def _poll_loop(self):
        last_offset = 0
        while self.is_running:
            try:
                resp = requests.get(
                    f"{self.base_url}/getUpdates",
                    params={"offset": last_offset, "timeout": 20},
                    timeout=25
                ).json()
                for update in resp.get("result", []):
                    last_offset = update["update_id"] + 1
                    msg = update.get("message", {})
                    chat_id = str(msg.get("chat", {}).get("id", ""))
                    if chat_id != self.paired_chat_id:
                        continue  # Reject unauthorized senders

                    text = msg.get("text", "")
                    if text == "/screenshot":
                        self._send_screenshot(chat_id)
                    elif text:
                        res = self.executor.execute_command(text)
                        self._send_message(chat_id, f"Executed: {res}")
            except Exception as e:
                time.sleep(5)

    def _send_message(self, chat_id: str, text: str):
        requests.post(f"{self.base_url}/sendMessage", json={"chat_id": chat_id, "text": text})

    def _send_screenshot(self, chat_id: str):
        try:
            import pyautogui, io
            img = pyautogui.screenshot()
            buf = io.BytesIO()
            img.save(buf, format="PNG")
            buf.seek(0)
            requests.post(f"{self.base_url}/sendPhoto", data={"chat_id": chat_id}, files={"photo": buf})
        except Exception as e:
            self._send_message(chat_id, f"Screenshot capture error: {e}")
\end{lstlisting}

Together, these five modules constitute the core operational engine of PRIVACY68, achieving complete hardware isolation, intelligent multi-provider reasoning, and deterministic security.
